\subsection{Compact dirty controls for selected-bit additions}
\label{sec:compact-control-movement}
The following replacement for the wide-slot packed operation moves only
compact control fields. Temporary address fields have arbitrary contents;
they are complete ranges of existing address bits, not newly allocated
independent coordinates. Their availability in every recursive call is
proved separately below.

\begin{lemma}[Controlled rotations with record-paid offset work]
\label{lem:record-paid-rotation}
Consider $M$ records of $R$ bits on a complete rectangular address set, of
volume $V=MR$. Its fixed number of field lengths and controls have at most
$A=\lceil\log_2(2V)\rceil$ bits each. A binary target field follows all its
controls. Suppose its rotation offset, independent of that target, is
computable from the controls and descriptors in $O(A^C)$ tape steps for
some fixed $C$. Then the rotation can be performed on fixed tapes in
$O(V+MA^{C'})$ time for a fixed $C'$. In particular, this is $O(V)$ when
$A=O(p)$ and $R$ exceeds every fixed power of $p$.
\end{lemma}
\begin{proof}
Fix the prefix before the target. If its range is $Q$ and each target block
has $B$ suffix bits, rotate by $a$ through the two-piece split at $Q-a$ and
merge in the opposite piece order. This is the controlled-rotation
construction from the ordered-affine streaming lemma. Each data bit is
copied a bounded number of times. There are at most $M$ prefix fibers, since
every fiber contains at least one complete record. At each fiber, compute
the offset, initialize its two piece lengths, and reset the finite collection
of counters; recompute the offset at the merge. Ordinary binary arithmetic
and descriptor copies cost a fixed polynomial in $A$ per fiber. The controls
can be wider than the target: their lengths are bounded by $A$, not by the
target width. Local suffix-length copies and ripple counters are charged as
in the original streaming proof, so no growing descriptor is scanned per
payload bit. Two piece tapes, output and fixed arithmetic work tapes suffice;
rewinds and clearing are over the regions just written. This gives the stated
bound. Finally $MA^{C'}/V=A^{C'}/R=o(1)$ in the specified record regime.
\end{proof}

Let $x,y$ be disjoint slots of $fK$ bits, with corresponding selected
positions $j_i=\rho+iK$, $0\le i<f$, $0\le\rho<K$. Put $n=f-1$.
For $f=1$ use the elementary two-bit XOR. Otherwise fix a positive integer
$G$, put $B=2^G$, and assume two disjoint complete $nG$-bit fields $u,t$
precede both slots and a third such field $b$ follows both slots. Arbitrary
spectator fields may intervene. Each compact field has $n$ radix-$B$ digits.
The selected-bit operation will preserve these three entire fields.

\paragraph{Earlier source.}
Suppose $x$ precedes $y$. The unused field $u$ is a spectator. For an integer
target segment $v$, a bit $z$, and an arbitrary integer temporary $w$, perform
\[
 v\gets v+2zw,\qquad w\gets w+(v\bmod2),\qquad
 v\gets v+z(1-2w),\qquad w\gets w-((v\bmod2)\mathbin\oplus z).
\]
Writing the original parity as $a$, the first update leaves it unchanged.
The final values are $v+z(1-2a)$ and the original $w$. Thus this toggles
only the parity of $v$ when $z=1$, and restores the dirty temporary.

Apply these four updates simultaneously, with $z_i=x_{j_i}$, to the low
$n$ selected segments of $y$ and the digits $t_i$ of $t$. A target-slot
update is one modular rotation of $y$, whose offsets are respectively
\[
 \sum_{i<n}2x_{j_i}t_i2^{j_i},\qquad
 \sum_{i<n}x_{j_i}(1-2t_i)2^{j_i} \pmod{2^{fK}}.
\]
Both control fields precede $y$. For a temporary update, swap $t$ with $b$,
rotate the back field by the packed selected bits of the current $y$, and
swap back. For unloading, use minus the packed bits
$y_{j_i}\mathbin\oplus x_{j_i}$. These controls precede the back field.
There are four width-$nG$ swaps and four rotations. The arbitrary original
contents of $b$ return after each pair of swaps.

\paragraph{Later source.}
If $x$ follows $y$, first execute the earlier-source construction with control
bits $u_i\bmod2$. Load the selected bits of $x$ into $u$ by swapping $u,b$,
rotating the back field by their packed radix-$B$ representation, and swapping
back. Repeat the earlier-source construction with the new parities of $u$,
then unload the selected bits of $x$ from $u$ by the inverse load.
When the digit loads do not overflow, the two control parities differ by
exactly $x_{j_i}$. The resulting target update is the required XOR, with
$u,t,b$ restored. This version uses twelve compact swaps and ten rotations.
Every offset depends on a fixed number of fields preceding its target.

\paragraph{Guards and deterministic repair.}
Temporarily omit the highest selected position. The low $n$ segments
$[j_i,j_i+K)$ are disjoint and fit inside $y$, including when $\rho=K-1$.
Write each initial segment as $2g_i+a_i$, $a_i\in\{0,1\}$.
Declare an address bad if a $t$ digit is $B-1$, or a target guard fails
\[
 2B\le g_i<2^{K-1}-2B.
\]
For a later source also declare it bad if any $u$ digit is $B-1$.
On a good address each temporary load stays within its digit, and the target
displacement is at most $2B$ in either execution of the four-step identity.
The conservative guard interval therefore prevents every carry or borrow
between the wide segments. After each completed identity the guard is
unchanged. Induction through the later-source sequence proves the desired
map on the good set, including restoration of every temporary and every
unselected address bit.

Uniform complete ranges give the common bound
\[
 \delta=\min\{1,n(2\cdot2^{-G}+8\cdot2^{G-K})\}
\]
on the bad fraction; for an earlier source the first factor 2 can be omitted.
The ideal map $T$ changes only the selected target parities, so it preserves
the bad set. The actual sequence $S$ is a bijection on the entire rectangle:
each rotation is independent of its target and all swaps are bijections.
Since $S=T$ off the bad set, $S$ preserves that set too. Thus $TS^{-1}$ is
supported there, exactly as in the original exceptional-stream repair.

Extract the bad records, retaining marked holes in the full stream. For each
extracted address $q$, reverse the swaps and modular rotations to compute
$S^{-1}q$, then compute $T(S^{-1}q)$. All fields and descriptors have $O(A)$
bits. A constant number of offset calculations each scan at most $n\le A$
digits and combine shifted integers of $O(A)$ bits. Integer arithmetic,
mixed-radix decoding, inverse offsets and the destination key can consequently
be computed in $O(A^3)$ per record. Sort only the extracted records by their
$O(A)$-bit destination keys, using stable binary radix passes, then reinsert
them in their marked holes. Keys are distinct and fill precisely those holes
because $TS^{-1}$ permutes the bad set. This uses fixed tapes and costs
\[
 O(V+MA^3+\delta MA(R+A)).
\]
Every inverse offset is evaluated on the current inverse controls; modular
overflow on bad addresses is handled exactly, rather than treated as an
integer identity. Finally, supply the omitted highest selected position
through the elementary two-bit XOR in $O(V)$ time.

\begin{proposition}[Compact-control selected-bit addition]
\label{prop:compact-selected-addition}
Under the stated complete-field hypotheses, take
\[
 G=4\lceil\log_2 p\rceil+6,\quad f\le p,\quad A=O(p),
 \quad K/\log p\longrightarrow\infty,
\]
and let $R$ exceed every fixed power of $p$. The exact selected-bit row
addition, with all temporary fields restored, costs
\[
 O\bigl(V((f\log p)^\tau+1)\bigr)
\]
on a fixed number of tapes, for either order of the two slots. Constants and
cutoffs are uniform within each fixed parameter family of the layer theorem.
\end{proposition}
\begin{proof}
Use the retained arbitrary-width chunk-swap primitive for the at most twelve
swaps, each of width $nG$. Its padding, cleanup and all spectator intervals
are covered by that interface. For every rotation, apply
Lemma~\ref{lem:record-paid-rotation}; its offset arithmetic is polynomial
in $A$. No new record volume is introduced by describing the compact fields.
For $\ell_p=\lceil\log_2 p\rceil$, eventually
$K\ge G+4\ell_p+10$. Then
\[
 \delta\le \frac{5}{128p^3},\qquad
 \frac{A^3}{R}=o(1),\qquad
 \delta A(1+A/R)=O(p^{-2}).
\]
The repair is therefore $O(V)$ per call, including predicate scans and key
arithmetic. The elementary highest-bit operation costs $O(V)$. All calls
reuse their fixed work tapes and finish in the declared physical field order.
\end{proof}

\paragraph{Application to a fixed binary change of basis.}
The original phase-frame construction expresses each binary basis change
as a fixed number of row additions on $m$ slots of $f=e/m$ axes. Use the
two slots of each addition as $x,y$ and the reserved front/back compact
fields as controls. The complete basis change and its inverse cost
$O(V((e\log p)^\tau+1))$. The reservations below guarantee their complete
ranges, independently of which slots are involved, at every recursive node.
